% Generated draft: requires mathematical and source review.
\documentclass[11pt]{article}
\usepackage[margin=25mm]{geometry}
\usepackage{fontspec}
\setmainfont{Latin Modern Roman}
\usepackage{amsmath,amssymb,mathtools,unicode-math}
\setmathfont{Latin Modern Math}
\usepackage{enumitem,needspace,xurl,hyperref}
\setlength{\parindent}{0pt}
\setlength{\parskip}{5pt}
\setlength{\emergencystretch}{3em}
\newcommand{\N}{\mathbb{N}}
\newcommand{\Z}{\mathbb{Z}}
\newcommand{\Q}{\mathbb{Q}}
\newcommand{\R}{\mathbb{R}}
\newcommand{\C}{\mathbb{C}}
\newcommand{\F}{\mathbb{F}}
\newcommand{\A}{\mathbb{A}}
\newcommand{\PP}{\mathbb{P}}
\newcommand{\E}{\mathbb{E}}
\newcommand{\eps}{\varepsilon}
\newcommand{\dd}{\,\mathrm{d}}
\begin{document}
\section*{UnsolvedMath ID: 30000699}
\textbf{Source problem code:} catalogue entry
\subsection{Definitions and mathematical statement}
Let \(\mathcal E\) denote the nonconstant entire functions \(\C\to\C\). For a finite set \(S\subset\C\) and \(f\in\mathcal E\), define the locally finite effective divisor
\[
E(S,f)=\sum_{a\in S}\operatorname{div}_0(f-a),
\]
where a zero of \(f-a\) is counted with its order of vanishing. Thus \(E(S,f)\) records the preimage of the whole set with multiplicities, without recording which value of \(S\) occurs at each point. If \(P_S(t)=\prod_{a\in S}(t-a)\), it is the zero divisor of \(P_S\circ f\).

Call \(S\) a finite unique-range set for nonconstant entire functions with multiplicities if
\[
\forall f,g\in\mathcal E,\qquad
E(S,f)=E(S,g)\ \Longrightarrow\ f\equiv g.
\]
The source asks to determine
\[
N_{\mathrm{ent}}=\min\{|S|:S\subset\C\text{ finite and has this property}\}.
\]
The set \(S\) is chosen before the universal pair \(f,g\). The question seeks one suitable set of the smallest size, not a claim that every set of that size is suitable. Identifying the minimum requires both an example at that size and exclusion of every smaller finite set. Set sharing is not replaced by separate labeled equalities of the divisors for each individual value, and multiplicity is not ignored. The functions are entire, not arbitrary meromorphic functions, and the value set contains no pole value at infinity.

The supplied literature assessment checked 2026-08-21 reports that existence of finite unique-range sets is known, while their optimal cardinality remains a quantitative problem with bounds and constructions rather than a resolved minimum. No numerical bound or extremal construction is reproduced, so none is invented.
\subsection{Short English statement}
Find the smallest size of one finite complex value set that determines every nonconstant entire function from its full preimage multiset. The value labels are forgotten but zero multiplicities are retained; existence alone does not determine the minimum.
\subsection{Sources}
\begin{itemize}
\item Local UnsolvedMath registry, permanent ID 30000699, catalogue entry, \emph{Minimal Unique-Range Sets for Entire Functions}.
\item \emph{Normal Families and Complex Dynamics} (2007), Oberwolfach Reports, pp. 487--548, DOI \url{https://doi.org/10.4171/owr/2007/09}.
\item Supplied reference ``Gross uniqueness question'': \url{https://www.cambridge.org/core/journals/bulletin-of-the-australian-mathematical-society/article/on-a-question-of-gross-concerning-uniqueness-of-entire-functions/4E1A936BE1AD28565E5F9EE7E16AD8DD}.
\item Supplied reference ``finite-set uniqueness theorems'': \url{https://doi.org/10.1080/17476930108815332}.
\item Supplied literature triage checked 2026-08-21 distinguishes existence from optimality; exact numerical bounds and construction hypotheses are absent.
\end{itemize}
\end{document}
